\chapter{Disjoint sets and spanning trees}\label{ch:dsu}
\index{disjoint set union}\index{DSU}\index{minimum spanning tree}

Two questions dominate this chapter: are these vertices connected,
and which edges are worth keeping? The disjoint-set union answers the
first in nearly constant time per operation; the minimum spanning
tree answers the second, and the cut property is the single theorem
that makes every MST algorithm correct. After this chapter you will
write a DSU blindfolded, argue Kruskal and Prim from the same
exchange proof, and spot when an offline replay turns ``delete''
into ``insert'' --- the trick that makes deletion queries solvable at
all.

\section{The connectivity problem}\label{sec:conn}\index{connectivity}
Maintain ``are $a$ and $b$ in the same component?'' under edge
insertions. Online algorithms must answer as each query arrives;
offline algorithms see the whole sequence first --- and offline is
where hard problems soften: deleting an edge sequence replayed
backwards is insertion, and DSU handles insertion natively.

\section{The forest representation}\label{sec:forest}
Each set is a rooted tree; \texttt{p[v]} points up. \texttt{find}
walks to the root, \emph{path compression} re-roots the walk's
vertices at the root; \texttt{unite} hangs the smaller root under the
larger (\emph{union by size}). Each half is optional; together they
are nearly free.
\begin{figure}[ht]\centering
% fig40-compress.tex -- Figure 8.1: find(d) with path compression.
% fstyle.tex -- shared TikZ styles for the figures of this book (palette from lifeos.sty).
\tikzset{%
  cell/.style={draw=ink, minimum width=7mm, minimum height=6.5mm, inner sep=0pt, font=\footnotesize, fill=white},
  cellh/.style={cell, fill=accentsoft, draw=accent},
  cella/.style={cell, fill=amberbg, draw=amber},
  cellg/.style={cell, fill=goodbg, draw=good},
  cellx/.style={cell, fill=graybg, draw=slate, text=slate},
  vtx/.style={circle, draw=ink, minimum size=6.5mm, inner sep=0pt, font=\footnotesize, fill=white},
  vtxh/.style={vtx, draw=accent, fill=accentsoft},
  vtxa/.style={vtx, draw=amber, fill=amberbg},
  vtxg/.style={vtx, draw=good, fill=goodbg},
  vtxd/.style={vtx, draw=slate, fill=graybg, text=slate},
  lab/.style={font=\footnotesize, text=slate},
  labA/.style={font=\footnotesize, text=accent},
  labM/.style={font=\footnotesize, text=amber},
  idx/.style={font=\tiny, text=slate},
  arr/.style={-{stealth}, thick, ink},
  arrA/.style={-{stealth}, thick, accent},
  arrM/.style={-{stealth}, thick, amber},
  arrG/.style={-{stealth}, thick, good},
  edge/.style={thick, ink},
  edgeA/.style={thick, accent},
  edgeM/.style={thick, amber},
}%
% \cellrow{prefix}{x0}{y}{style}{v0,v1,...}: a row of cells named prefix0, prefix1, ... (pitch 0.75cm)
\providecommand{\cellrow}[5]{\foreach \v [count=\i from 0] in {#5} {\node[#4] (#1\i) at (#2+\i*0.75,#3) {\v};}}%
% \cellidx{prefix}{n}{start}: index labels under cells prefix0 .. prefix(n-1)
\providecommand{\cellidx}[3]{\foreach \i in {0,...,\numexpr#2-1\relax} {\node[idx, below=0pt of #1\i] {\the\numexpr\i+#3\relax};}}%
% \fitfig{<tikzpicture>}: shrink the picture to \linewidth only when it would overflow.
\ifdefined\fitbox\else\newsavebox{\fitbox}\fi%
\providecommand{\fitfig}[1]{\sbox{\fitbox}{#1}\ifdim\wd\fitbox>\linewidth\resizebox{\linewidth}{!}{\usebox{\fitbox}}\else\usebox{\fitbox}\fi}%

\begin{tikzpicture}[font=\footnotesize]
% before
\node[vtx] (r) at (0,3.0) {r}; \node[vtx] (a) at (0,2.1) {a}; \node[vtxh] (b) at (0,1.2) {b}; \node[vtxh] (c) at (0,0.3) {c};
\node[vtxh] (d) at (0,-0.6) {d}; \node[vtx] (x) at (1.1,0.3) {x};
\draw[arr, shorten >=1pt] (a)--(r); \draw[arr, shorten >=1pt] (b)--(a); \draw[arr, shorten >=1pt] (c)--(b); \draw[arr, shorten >=1pt] (d)--(c); \draw[arr, shorten >=1pt] (x)--(b);
\node[lab, align=center] at (0.5,-1.5) {before \texttt{find(d)}};
\draw[arrA, ultra thick] (2.2,1.2) -- (3.2,1.2) node[midway, above, labA] {find};
% after
\node[vtx] (r2) at (5.5,3.0) {r};
\node[vtx] (a2) at (4.3,1.7) {a}; \node[vtxh] (b2) at (5.1,1.7) {b}; \node[vtxh] (c2) at (5.9,1.7) {c}; \node[vtxh] (d2) at (6.7,1.7) {d};
\node[vtx] (x2) at (5.1,0.3) {x};
\foreach \n in {a2,b2,c2,d2} \draw[arr, shorten >=1pt] (\n)--(r2);
\draw[arr, shorten >=1pt] (x2)--(b2);
\node[lab, align=center] at (5.5,-0.55) {after: path re-parented to the root,\\ off-path $x$ untouched};
\end{tikzpicture}

\caption{The same find, before and after. Notice compression only
re-parents vertices on the walked path --- others keep their depth
until their turn.}
\label{fig:compress}
\end{figure}
\begin{figure}[ht]\centering
\begin{tikzpicture}[font=\footnotesize]
\begin{scope}
\node[circle,draw=slate,fill=accentsoft,font=\tiny] at (0,1) {big};
\node[circle,draw=slate,font=\tiny] at (0.8,1) {small};
\draw[-{stealth},amber,thick] (0.8,1) to[out=45,in=45] (0,1);
\node[font=\tiny] at (0.4,0.4) {small $\to$ big: stays shallow};
\end{scope}
\begin{scope}[xshift=4.4cm]
\node[circle,draw=slate,font=\tiny] at (0,1) {big};
\node[circle,draw=slate,fill=accentsoft,font=\tiny] at (0.8,1) {small};
\draw[-{stealth},amber,thick,dashed] (0,1) to[out=135,in=135] (0.8,1);
\node[font=\tiny] at (0.4,0.4) {big $\to$ small: deepens};
\end{scope}
\end{tikzpicture}

\caption{Union by size. Notice the asymmetry: only hanging the larger
under the smaller deepens anything; the other direction never adds
depth.}
\label{fig:unionsize}
\end{figure}
\begin{lstlisting}[language=C++,caption={\texttt{dsu.h}: path
halving + union by size + component sizes.},label={lst:dsu}]
\end{lstlisting}
\lstinputlisting{code/dsu.h}

\section{The amortized bound}\label{sec:ack}
With both optimizations, $m$ operations cost $O(m\,\alpha(n))$ where
$\alpha$ is the inverse Ackermann function --- bounded by 4 for any
$n$ in this universe. The proof charges each pointer hop to a rank
increase budget; sketch only: ranks are potential, compression spends
it. In contest terms: treat each operation as $O(1)$.

\section{The cut property}\label{sec:cut}\index{cut property}
For any cut $(S, V\setminus S)$, the lightest crossing edge belongs
to \emph{some} MST. Proof by exchange: take an MST using a heavier
crossing edge $f$; swapping in the light edge $e$ keeps a spanning
tree and only lowers the cost. Every MST algorithm is this theorem in
a different costume.
\begin{figure}[ht]\centering
% fig42-cutprop.tex -- Figure 8.3: the cut property; the lightest edge crossing a cut belongs to some MST.
% fstyle.tex -- shared TikZ styles for the figures of this book (palette from lifeos.sty).
\tikzset{%
  cell/.style={draw=ink, minimum width=7mm, minimum height=6.5mm, inner sep=0pt, font=\footnotesize, fill=white},
  cellh/.style={cell, fill=accentsoft, draw=accent},
  cella/.style={cell, fill=amberbg, draw=amber},
  cellg/.style={cell, fill=goodbg, draw=good},
  cellx/.style={cell, fill=graybg, draw=slate, text=slate},
  vtx/.style={circle, draw=ink, minimum size=6.5mm, inner sep=0pt, font=\footnotesize, fill=white},
  vtxh/.style={vtx, draw=accent, fill=accentsoft},
  vtxa/.style={vtx, draw=amber, fill=amberbg},
  vtxg/.style={vtx, draw=good, fill=goodbg},
  vtxd/.style={vtx, draw=slate, fill=graybg, text=slate},
  lab/.style={font=\footnotesize, text=slate},
  labA/.style={font=\footnotesize, text=accent},
  labM/.style={font=\footnotesize, text=amber},
  idx/.style={font=\tiny, text=slate},
  arr/.style={-{stealth}, thick, ink},
  arrA/.style={-{stealth}, thick, accent},
  arrM/.style={-{stealth}, thick, amber},
  arrG/.style={-{stealth}, thick, good},
  edge/.style={thick, ink},
  edgeA/.style={thick, accent},
  edgeM/.style={thick, amber},
}%
% \cellrow{prefix}{x0}{y}{style}{v0,v1,...}: a row of cells named prefix0, prefix1, ... (pitch 0.75cm)
\providecommand{\cellrow}[5]{\foreach \v [count=\i from 0] in {#5} {\node[#4] (#1\i) at (#2+\i*0.75,#3) {\v};}}%
% \cellidx{prefix}{n}{start}: index labels under cells prefix0 .. prefix(n-1)
\providecommand{\cellidx}[3]{\foreach \i in {0,...,\numexpr#2-1\relax} {\node[idx, below=0pt of #1\i] {\the\numexpr\i+#3\relax};}}%
% \fitfig{<tikzpicture>}: shrink the picture to \linewidth only when it would overflow.
\ifdefined\fitbox\else\newsavebox{\fitbox}\fi%
\providecommand{\fitfig}[1]{\sbox{\fitbox}{#1}\ifdim\wd\fitbox>\linewidth\resizebox{\linewidth}{!}{\usebox{\fitbox}}\else\usebox{\fitbox}\fi}%

\begin{tikzpicture}[font=\footnotesize]
\draw[rounded corners=8pt, draw=accent, fill=accentsoft] (-0.2,-0.2) rectangle (3.4,3.6);
\draw[rounded corners=8pt, draw=amber, fill=amberbg] (5.2,-0.2) rectangle (8.8,3.6);
\node[labA] at (1.6,3.9) {$S$}; \node[labM] at (7.0,3.9) {$V\setminus S$};
\node[vtx] (p1) at (0.8,2.8) {}; \node[vtx] (p2) at (2.4,2.6) {}; \node[vtx] (p3) at (1.2,0.9) {}; \node[vtx] (p4) at (2.6,0.6) {};
\node[vtx] (q1) at (6.0,2.7) {}; \node[vtx] (q2) at (7.8,2.9) {}; \node[vtx] (q3) at (6.2,0.8) {}; \node[vtx] (q4) at (7.9,0.8) {};
\draw[edge, slate] (p1)--(p2) (p1)--(p3) (p3)--(p4) (q1)--(q2) (q1)--(q3) (q3)--(q4);
\draw[edge, slate] (p2) -- (q1) node[midway, above, lab] {9};
\draw[edgeM, very thick] (p4) -- (q3) node[midway, below, labM] {3};
\draw[edge, slate] (p2) -- (q4) node[pos=0.25, below left=1pt, lab] {6};
\node[labM, align=center] at (4.0,-1.0) {lightest crossing edge (3): safe to take};
\draw[dashed, slate] (4.0,-0.4) -- (4.0,3.9);
\end{tikzpicture}

\caption{The cut property. Notice nothing depends on which MST you
started from --- exchange against \emph{any} of them.}
\label{fig:cutprop}
\end{figure}

\section{Kruskal}\label{sec:kruskal}\index{Kruskal}
Sort edges, accept those joining different components. Each accepted
edge is the lightest crossing of the cut between the two components
--- the cut property again, applied greedily. Cost: sorting
dominates, $O(E\log E)$.
\begin{lstlisting}[language=C++,caption={\texttt{kruskal.h} on top of
\texttt{DSU}.},label={lst:kruskal}]
\end{lstlisting}
\lstinputlisting{code/kruskal.h}
\begin{example}[CSES \emph{Road Reparation}]\label{ex:reparation}
Cheapest way to reconnect the cities.
\par\smallskip
\textbf{Strategy.} ``Cheapest connection covering all'' is an MST;
a leftover of fewer than $n{-}1$ accepted edges means the graph was
never connected.
\par\smallskip
\begin{solution}
\lstinputlisting{code/road_reparation.cpp}
Real output (CSES sample):
\begin{lstlisting}[numbers=none]
14
\end{lstlisting}
\end{solution}
\end{example}
\begin{figure}[ht]\centering
% fig43-kruskal.tex -- Figure 8.4: Kruskal on the CSES Road Reparation sample (answer 14).
% fstyle.tex -- shared TikZ styles for the figures of this book (palette from lifeos.sty).
\tikzset{%
  cell/.style={draw=ink, minimum width=7mm, minimum height=6.5mm, inner sep=0pt, font=\footnotesize, fill=white},
  cellh/.style={cell, fill=accentsoft, draw=accent},
  cella/.style={cell, fill=amberbg, draw=amber},
  cellg/.style={cell, fill=goodbg, draw=good},
  cellx/.style={cell, fill=graybg, draw=slate, text=slate},
  vtx/.style={circle, draw=ink, minimum size=6.5mm, inner sep=0pt, font=\footnotesize, fill=white},
  vtxh/.style={vtx, draw=accent, fill=accentsoft},
  vtxa/.style={vtx, draw=amber, fill=amberbg},
  vtxg/.style={vtx, draw=good, fill=goodbg},
  vtxd/.style={vtx, draw=slate, fill=graybg, text=slate},
  lab/.style={font=\footnotesize, text=slate},
  labA/.style={font=\footnotesize, text=accent},
  labM/.style={font=\footnotesize, text=amber},
  idx/.style={font=\tiny, text=slate},
  arr/.style={-{stealth}, thick, ink},
  arrA/.style={-{stealth}, thick, accent},
  arrM/.style={-{stealth}, thick, amber},
  arrG/.style={-{stealth}, thick, good},
  edge/.style={thick, ink},
  edgeA/.style={thick, accent},
  edgeM/.style={thick, amber},
}%
% \cellrow{prefix}{x0}{y}{style}{v0,v1,...}: a row of cells named prefix0, prefix1, ... (pitch 0.75cm)
\providecommand{\cellrow}[5]{\foreach \v [count=\i from 0] in {#5} {\node[#4] (#1\i) at (#2+\i*0.75,#3) {\v};}}%
% \cellidx{prefix}{n}{start}: index labels under cells prefix0 .. prefix(n-1)
\providecommand{\cellidx}[3]{\foreach \i in {0,...,\numexpr#2-1\relax} {\node[idx, below=0pt of #1\i] {\the\numexpr\i+#3\relax};}}%
% \fitfig{<tikzpicture>}: shrink the picture to \linewidth only when it would overflow.
\ifdefined\fitbox\else\newsavebox{\fitbox}\fi%
\providecommand{\fitfig}[1]{\sbox{\fitbox}{#1}\ifdim\wd\fitbox>\linewidth\resizebox{\linewidth}{!}{\usebox{\fitbox}}\else\usebox{\fitbox}\fi}%

\begin{tikzpicture}[font=\footnotesize]
\node[vtx] (n1) at (0,2) {1}; \node[vtx] (n2) at (2.2,2) {2}; \node[vtx] (n3) at (4.4,2) {3};
\node[vtx] (n5) at (0,0) {5}; \node[vtx] (n4) at (2.2,0) {4};
\draw[edge, slate, dashed] (n1)--(n5) node[midway, left, lab] {7};
\draw[edge, slate, dashed] (n3)--(n4) node[midway, below right, lab] {8};
\draw[edgeM, very thick] (n1)--(n2) node[midway, above, labM] {3};
\draw[edgeM, very thick] (n2)--(n3) node[midway, above, labM] {5};
\draw[edgeM, very thick] (n2)--(n4) node[midway, right, labM] {2};
\draw[edgeM, very thick] (n5)--(n4) node[midway, below, labM] {4};
\node[lab, align=left, anchor=west] at (5.6,2.35) {edges sorted by weight:};
\foreach \y/\w/\e/\verdict in {1.7/2/{2--4}/{accept}, 1.2/3/{1--2}/{accept}, 0.7/4/{4--5}/{accept}, 0.2/5/{2--3}/{accept}, -0.3/7/{1--5}/{skip (same set)}, -0.8/8/{3--4}/{skip (same set)}}
  \node[lab, anchor=west] at (5.6,\y) {$\w$: edge \e\ $\to$ \verdict};
\node[labM] at (7.4,-1.5) {total $2+3+4+5=14$};
\end{tikzpicture}

\caption{Kruskal's accept/skip rhythm. Notice the skipped edge is
rejected \emph{because} its endpoints already merged --- the DSU is
the algorithm's memory.}
\label{fig:kruskal}
\end{figure}

\section{Prim}\label{sec:prim}\index{Prim}
Grow one tree: a priority queue holds frontier edges; pop the
lightest, settle its far vertex, push its edges. It is Dijkstra's
skeleton with edge weight instead of distance --- and wins when the
graph is dense enough that sorting all edges is the bottleneck.
\begin{lstlisting}[language=C++,caption={\texttt{prim.h}: the pq
frontier.},label={lst:prim}]
\end{lstlisting}
\lstinputlisting{code/prim.h}

\section{Worked examples}
\begin{example}[CSES \emph{Road Construction}]\label{ex:constr}
After each new road: component count, largest component.
\par\smallskip
\textbf{Strategy.} Both answers live in the DSU's size array:
merges decrement the count and grow one component.
\par\smallskip
\begin{solution}
\lstinputlisting{code/road_construction.cpp}
Real output (CSES sample):
\begin{lstlisting}[numbers=none]
4 2
3 3
2 3
\end{lstlisting}
\end{solution}
\end{example}
\begin{example}[CSES \emph{Building Roads}]\label{ex:roads}
Minimum roads to connect all cities, then print them.
\par\smallskip
\textbf{Strategy.} Components minus one is the count; linking
component roots in a chain realizes it.
\par\smallskip
\begin{solution}
\lstinputlisting{code/building_roads.cpp}
Real output (CSES sample):
\begin{lstlisting}[numbers=none]
1
1 3
\end{lstlisting}
\end{solution}
\end{example}
\begin{example}[USACO-S \emph{Closing the Farm}]\label{ex:farm}
After each barn closes, is the farm still fully connected?
\par\smallskip
\textbf{Strategy.} DSU cannot delete --- so run time backwards:
start from the last surviving barn and \emph{insert} barns in reverse
order, recording the connectivity of each state \emph{before}
the closed barn is re-inserted (the answer asks for the state
\emph{after} that closing; an empty remainder is vacuously
connected)
(Figure~\ref{fig:farmtime}).
\par\smallskip
\begin{solution}
\lstinputlisting{code/closing_farm.cpp}
Real output (4-barn star, all states connected):
\begin{lstlisting}[numbers=none]
YES
YES
YES
YES
\end{lstlisting}
\end{solution}
\end{example}
\begin{figure}[ht]\centering
\begin{tikzpicture}[font=\footnotesize]
\node[font=\tiny] at (0,0.5) {delete 4, 3, 1, 2};
\draw[-{stealth},slate] (0,0.2)--(2.6,0.2);
\node[font=\tiny] at (3.4,0.2) {time};
\begin{scope}[yshift=-0.7cm]
\node[font=\tiny] at (0,0.3) {add 2, 1, 3, 4};
\draw[-{stealth},amber,thick] (2.6,-0.1)--(0,-0.1);
\node[font=\tiny] at (3.4,0) {replay};
\end{scope}
\node[font=\tiny,align=center] at (1.3,-1.4) {offline: answer each prefix\\after re-adding in reverse};
\end{tikzpicture}

\caption{The timeline reversal. Notice the forward question and the
backward algorithm share \emph{the same states} --- replaying
backwards visits every intermediate farm the original run did.}
\label{fig:farmtime}
\end{figure}

\begin{exercise}[easy]\label{exr:fu}
In \texttt{unite(a,b)}, which order runs \texttt{find} first --- and
why does union-by-size need the \emph{root} sizes, not arbitrary
sizes?
\end{exercise}
\begin{exercise}[medium]\label{exr:whichcut}
Kruskal's next accepted edge $e$ between components $A,B$: exhibit
the cut whose light edge $e$ is, and conclude its membership.
\end{exercise}
\begin{exercise}[medium]\label{exr:reverse}
Give a query sequence where offline reversal is essential (not just
convenient): name one operation DSU fundamentally cannot undo.
\end{exercise}

\section{Chapter summary}
\begin{table}[ht]\centering
{\small
\begin{tabularx}{\linewidth}{@{}lYYY@{}}
\toprule
\tablehead\thead{Task} & \thead{Tool} & \thead{Cost} & \thead{Watch for}\\
\midrule
connected? + merge & DSU & $\sim O(1)$/op & no deletes\\
all-pairs min tree & Kruskal & $O(E\log E)$ & sort dominates\\
dense-graph MST & Prim & $O(E\log V)$ & pq frontier\\
deletion queries & offline reverse & same ops & replay order\\
\bottomrule
\end{tabularx}}
\caption{One theorem (cut property) under both MST algorithms; one
trick (offline reversal) under every deletion problem.}
\label{tab:dsu}
\end{table}

\section{Common mistakes}
\index{common mistakes}
\begin{itemize}
\item \textbf{Union without \texttt{find}.} \texttt{p[a]=p[b]} links
arbitrary members, not roots --- always unite
\texttt{find(a)},\,\texttt{find(b)}.
\item \textbf{Skipping compression \emph{and} size.} One alone helps;
neither degenerates to $O(n)$ per find --- keep both.
\item \textbf{Trying to delete with DSU.} There is no
\texttt{un-unite}; offline reversal or a different structure is the
answer.
\item \textbf{Assuming MST is unique.} Equal weights admit several;
the cut property promises \emph{some} MST, not \emph{the} MST ---
handle ties deliberately when the problem counts trees.
\item \textbf{Forgetting the disconnected case.} A forest is legal
input --- \texttt{mst.size() < n-1} is the IMPOSSIBLE check.
\end{itemize}

\section{Additional exercises}
\begin{exercise}[easy] CSES \emph{Road Reparation}. \end{exercise}
\begin{exercise}[easy] USACO-S \emph{Closing the Farm}. \end{exercise}
\begin{exercise}[easy] USACO-G \emph{MooTube} --- threshold components. \end{exercise}
\begin{exercise}[medium] USACO-S \emph{Wormhole Sort} --- DSU + binary search. \end{exercise}
\begin{exercise}[medium] USACO-G \emph{Moocast} --- MST-style radii. \end{exercise}
\begin{exercise}[medium] USACO-G \emph{I Would Walk 500 Miles} --- implicit-edge Kruskal. \end{exercise}
\begin{exercise}[medium] CF \emph{Design Tutorial: Inverse the Problem} --- reconstruct from MST. \end{exercise}
\begin{exercise}[hard] CSES \emph{New Roads Queries} --- online MST edge questions. \end{exercise}
\begin{exercise}[hard] USACO-G \emph{Portals} --- MST on portal+cost edges. \end{exercise}
\begin{exercise}[hard] USACO-G \emph{Fenced In} --- boundary-strip Kruskal. \end{exercise}
\begin{exercise}[hard] USACO-OS \emph{Ski Course Rating} --- MST + queries. \end{exercise}
\begin{exercise}[hard] Baltic OI 2016 \emph{Park} --- MST + LCA path max. \end{exercise}

\chapternotes
DSU amortization and the cut-property-first MST treatment follow
Laaksonen~\cite{laaksonen2024} and Sannemo~\cite{sannemo2025}; the
offline-reversal pattern and threshold-component problems are the
USACO Guide's DSU module~\cite{usacoguide}; cp-algorithms carries the
charged-ranks proof sketch and MST variants~\cite{cpalgorithms};
Halim supplies the named variants~\cite{halimcp4}; archive problems
are official CSES/USACO~\cite{cses}.
